\scanpage{26}\mspage{24}
il faut prouver que les deux sont homotopes. On
a donc deux foncteurs
\diagram{\begin{tikzcd}[ampersand replacement=\&,column sep=1.5em]
(*)\quad f,g\colon
\underbrace{\overbrace{\underline{\mathrm{Fl}}(X;x,y)}^{\text{cat. discr.}}\times
\overbrace{\underline{\mathrm{Fl}}(X,y,z)}^{\text{cat. discr.}}}_{A}
\arrow[r,no head,shift left=1,"f"]\arrow[r,no head,shift right=1,"g"']
\& \underbrace{\unsure{\mathrm{Ch}|\mathscr{L}}_{4,\unsure{\Sigma|Z}}(X;x,z)}_{B}
\end{tikzcd}},
on veut prouver qu'ils sont homotopes. Pour ceci,
on les regarde \deleted{comme tantôt} \inserted{(comme tantôt)} comme ``transposés'' de deux
foncteurs
\diagram{\begin{tikzcd}[ampersand replacement=\&,column sep=large]
\mathord{\tikz[baseline=0pt,x=1ex,y=1ex]{\draw[line width=.09ex,line join=round] (0,0) -- (.95,1.30) -- (1.52,0) -- cycle; \draw[line width=.09ex] (.32,0) -- (1.02,1.12);}}_1\unsure{\amalg|\sqcup}\mathord{\tikz[baseline=0pt,x=1ex,y=1ex]{\draw[line width=.09ex,line join=round] (0,0) -- (.95,1.30) -- (1.52,0) -- cycle; \draw[line width=.09ex] (.32,0) -- (1.02,1.12);}}_1\simeq\mathord{\tikz[baseline=0pt,x=1ex,y=1ex]{\draw[line width=.09ex,line join=round] (0,0) -- (.95,1.30) -- (1.52,0) -- cycle; \draw[line width=.09ex] (.32,0) -- (1.02,1.12);}}_2
\& \unsure{\mathscr{L}|\mathcal{L}}_4 \arrow[l,shift right=1,"\varphi"']\arrow[l,shift left=1,"\psi"]\\
{} \& \varepsilon\arrow[u,"\alpha'"']\arrow[ul,"\alpha"']
\end{tikzcd}},
vu que les catégories \inserted{$A,B$} des deux \deleted{\unsure{buts}} de $(*)$
s'interprètent comme des sous-catégories pleines
de $\underline{\mathrm{Hom}}(\mathord{\tikz[baseline=0pt,x=1ex,y=1ex]{\draw[line width=.09ex,line join=round] (0,0) -- (.95,1.30) -- (1.52,0) -- cycle; \draw[line width=.09ex] (.32,0) -- (1.02,1.12);}}_2,X)$ et $\underline{\mathrm{Hom}}(\mathscr{L}_4,X)$ respectivement ;
à savoir des sous-catégories pleines formées
des foncteurs dont le composé avec $\alpha$ resp.$^{\mathrm{t}}$
$\alpha'$ est donné par $x$ et par $z$, et de plus
\deleted{\unsure{du fait de}}\unsure{(pour} $A$) dont la valeur en l'objet médian de
$\mathord{\tikz[baseline=0pt,x=1ex,y=1ex]{\draw[line width=.09ex,line join=round] (0,0) -- (.95,1.30) -- (1.52,0) -- cycle; \draw[line width=.09ex] (.32,0) -- (1.02,1.12);}}_2$ soit $y$. On peut encore remplacer $A$
par la ss-catégorie pleine plus grande \inserted{$A'$} de
$\underline{\mathrm{Hom}}(\mathord{\tikz[baseline=0pt,x=1ex,y=1ex]{\draw[line width=.09ex,line join=round] (0,0) -- (.95,1.30) -- (1.52,0) -- cycle; \draw[line width=.09ex] (.32,0) -- (1.02,1.12);}}_2,X)$, obtenue en ignorant cette
dernière condition, et il \deleted{\unsure{est évident}} \inserted{suffit de}
prouver que les deux foncteurs ``plus gros''
\diagram{\begin{tikzcd}[ampersand replacement=\&]
A'\arrow[r,shift left=1,"f'"]\arrow[r,shift right=1,"{g^{\unsure{\prime|=}}}"']\& B
\end{tikzcd}}
sont homotopes. Et pour ceci, il suffit de
% archive sheet number: 25
